\documentclass{article}

\usepackage{graphicx} % Required for inserting images
\usepackage{amsmath} % For \text in math mode
\usepackage[backend=biber]{biblatex}
\addbibresource{Bibliography.bib}

\title{32}
\author{c.marquezb }
\date{April 2026}

\begin{document}

\maketitle

\section{Introduction}

Ensuring successful execution of a program is a key requirement in software engineering. In classical computing, this process is done by directly comparing the observed output against an expected value. In other words, each execution produces a reproducible output. However, this deterministic validation model is not applicable in the QC realm, since in order to take advantage of quantum mechanics, it is required to run an algorithm multiple times and then analyze the resulting distribution to validate the execution success. In particular, a single run of a quantum algorithm produces a single classical state \cite{nielsenQuantumComputationQuantuma_2010}. Yet, to take advantage of quantum phenomena such as entanglement, superposition, and interference, it is necessary to construct a probability distribution over all possible outcomes by running the algorithm multiple times. Today, a histogram is typically used to represent this output, where the x-axis corresponds to the possible bit-string outcomes, and the y-axis counts the frequency of the observed values. By contrast, in hardware based on quantum annealing \cite{Hauke_2020_2020}, the output corresponds to the physical configuration of the system (a specific bit string) that represents the ground state (or a low-energy configuration). Another alternative to histograms is seen in variational algorithms \cite{McClean_2016_2016}, where quantum hardware outputs a real scalar value (the expected value of a cost function, such as energy), which is used as input to a classical optimizer to finally produce a bit-string result. The probability distributions described above are accessible through the SDKs or APIs provided by the quantum hardware vendors, most of which use superconducting qubits for the built-in quantum processing unit (QPU). A common feature of such SDKs is the ability to compute histograms of execution outcomes. This motivates the creation of a metric that quantifies how distinctly the expected outcomes stand out within the observed distribution.

Before introducing such a metric, it is worth considering the applicability of classical performance metrics in quantum software; although doing so restricts the evaluation of quantum success to a classical paradigm in which the observed outcomes are fundamentally stochastic. More precisely, classical success measures can be defined at the shot level, but, because quantum processes require repeated executions to obtain meaningful results, the success measure should be formulated at the job level and, in some scenarios, even at the batch level. Several metrics have been proposed to quantify the successful execution of quantum circuits \cite{mohapatraBenchmarkingFidelityMetrics_2025}; however, most of them are designed to characterize the fidelity of the quantum hardware rather than to evaluate how distinctly a target outcome emerges in the observed distribution.

% Section~\ref{sec:related-work} explores these metrics in detail and positions the DSR relative to them. In this paper, we introduce the differential success rate (DSR), a quantitative metric designed to measure the level of distinction of a target quantum state within a given histogram. The DSR is expressed as a percentage, thereby serving as an indicator of how clearly the expected quantum state emerges relative to the observed probability distribution generated by current noisy intermediate-scale quantum (NISQ) devices. To characterize the behavior of DSR, we first execute the Grover and quantum Fourier transform algorithms using the AER simulator parameterized with the hardware calibration values reported for QPUs commonly available from major quantum service providers. The DSR metric is then integrated into qWard, a Python-based framework for the analysis of quantum circuits. Finally, we run the corresponding circuits on the physical QPUs \textit{ibm_fez}, \textit{ibm_torino}, and \textit{ibm_marrakesh}---three instances of the IBM Heron~R2 processor---accessed via the IBM Quantum service, as well as on the Rigetti \textit{Ankaa-3} device accessed via Amazon Web Services (AWS).

Our experiments demonstrate the usefulness of the DSR metric for improving the interpretability of quantum jobs by defining clear ranges of success for the Grover and QFT algorithms, on two quantum providers (4 different QPUs). Such intervals are determined by systematically scaling and executing the algorithms, computing the DSR, and correlating its values with the corresponding number of qubits and the circuit depth. In addition, we empirically compare DSR against HF and TVDF, showing that the metrics capture complementary aspects of algorithm success.

The remainder of this paper is organized as follows. Section~\ref{sec:background} presents background work. Section~\ref{sec:related-work} reviews related fidelity metrics. Section~\ref{sec:measuring-success} introduces the execution model and success criteria, while Section~\ref{sec:dsr} defines the DSR metric and evaluates alternative formulations. Section~\ref{subsec:experiment} describes the setup of the experiment, including workloads, backends, and variables. Section~\ref{sec:results} reports the empirical results: the relationship between DSR and qubit count, circuit depth, optimization level, and a statistical comparison with Hellinger and TVD fidelity. Section~\ref{sec:discussion} discusses limitations and validity considerations, and Section~\ref{sec:conclusion} concludes the paper.

\section{Methods}

[Move: Overview and Research Aims] To characterize the behavior of the Differential Success Rate (DSR) and benchmark its performance, we evaluate the execution of two well-defined algorithms. Our primary aim is to systematically scale these algorithms, compute the DSR from the resulting histograms, and statistically correlate these values with circuit depth, qubit count, and hardware provider noise levels. In this section, we outline the workloads, the quantum hardware utilized, the experimental pipeline, and the statistical tools used for data analysis.

[Move: Subjects and Materials] To ensure a rigorous characterization, we select the Grover search algorithm and the Quantum Fourier Transform (QFT). These workloads operate on pre defined quantum states, which is a requirement to calculate the DSR score. The dataset comprises 1,281 quantum jobs: 246 for Grover's search, 239 for QFT, and 796 for a variation of the teleportation protocol (vTP) from a prior study. For the Grover algorithm, the configurations range from 2 to 14 qubits for the scalability study and incorporate variations in marked count, Hamming weight, and symmetry. Conversely, QFT circuits are evaluated using both a round-trip mode spanning 2 to 10 qubits and a period-detection mode spanning 3 to 6 qubits.

To execute these circuits, we employ superconducting quantum processors provided by two major cloud-based quantum computing platforms. In particular, we employ multiple IBM Heron R2 quantum processors (specifically, the devices fez, torino, kingston, and marrakesh), which are accessed remotely via the IBM Quantum cloud service. In addition, we used the Rigetti Ankaa-3 device through Amazon Web Services (AWS).

[Move: Procedure] The execution pipeline is organized into three consecutive phases: development, execution, and analysis. First, to verify the correctness of the code, configurations are simulated using the Qiskit AER simulator. Following the successful simulation, the algorithms are executed on real quantum processing units (QPUs). During this phase, 1,024 shots are allocated per job. To enable fair comparisons between providers, the circuits are transpiled using the default provider settings. In IBM backends, each configuration is transpiled at four distinct optimization levels, ranging from 0 through 3. Meanwhile, circuits executed on the Rigetti backend are transpiled strictly at optimization level 3, as the provider does not expose lower optimization tiers.

[Move: Data Analysis] Once the execution is complete, the histograms of the outcomes are collected, and post-runtime metrics are calculated. The dependent variable is the DSR, computed from the target state and the returned histogram. To contextualize the DSR with respect to established fidelity measures, we additionally compute the Hellinger fidelity (HF) and the total variation distance fidelity (TVDF) for each QPU execution.

To compare DSR against HF and TVDF, we applied three non-parametric analyzes. The differences compared within each qubit group were evaluated with the Wilcoxon signed-rank test~\cite{Wilcoxon1992_1992}, using Bonferroni-corrected $p$-values to account for the two comparisons per group (DSR vs.\ HF and DSR vs.\ TVDF)~\cite{McEwan_2017_2017}. We also report Cliff's~$\delta$~\cite{cliffDominanceStatisticsOrdinal_1993} as an effect size, and bootstrap 95\% confidence intervals based on $10{,}000$ samples for the median of each metric in groups with $n \geq 5$~\cite{efronIntroductionBootstrap_1994}.

[Move: Limitations] Despite the systematic nature of this experimental setup, the methodology is subject to certain limitations. Because DSR requires knowledge of the expected outcomes, this approach is not directly applicable to algorithms whose ideal output is a broad probability distribution without a discrete set of dominant states, such as variational algorithms. In such cases, distribution-level metrics, such as HF or TVDF, are more appropriate. Furthermore, the precision of the target quantum states~$E$ is essential for an accurate evaluation. A limited number of shots, particularly when acquired with noisy hardware, introduces sampling uncertainty that propagates into the DSR value. Finally, backend calibration variation between jobs may introduce variability that is attributable to hardware fluctuations rather than the algorithmic configuration itself.

\printbibliography

\end{document}
